% !TEX TS-program = pdflatex
% FIVE_DAY_WRAPUP_PLAN.tex
% Day-by-day plan to close the whole project: design closure, 4-tile bench
% prototype, test, thesis assembly, submission. Day 1 = Thu 11 June 2026,
% Day 5 = Mon 15 June 2026. Companion document: THESIS_OUTLINE.tex.
% Format follows PROJECT_COMPLETION_PLAN.tex.
\documentclass[11pt,a4paper]{article}

\usepackage[margin=2.3cm]{geometry}
\usepackage[T1]{fontenc}
\usepackage{lmodern}
\usepackage{microtype}
\usepackage{setspace}
\onehalfspacing
\usepackage{amsmath,amssymb}
\usepackage{booktabs}
\usepackage{longtable}
\usepackage{colortbl}
\usepackage[table]{xcolor}
% Compat shim: colortbl 1.0k (2024) expects \insert@pcolumn from array v2.6,
% but this MiKTeX tree carries array v2.5g. Without the alias every p-column
% in a colortbl document throws "Undefined control sequence".
\makeatletter
\@ifundefined{insert@pcolumn}{\let\insert@pcolumn\insert@column}{}
\makeatother
\usepackage{enumitem}
\usepackage[hidelinks]{hyperref}
\usepackage{titlesec}
\titleformat{\section}{\large\bfseries}{\thesection.}{0.5em}{}
\titleformat{\subsection}{\normalsize\bfseries}{\thesubsection}{0.5em}{}
\emergencystretch=1.5em   % long typewriter file names need slack

\definecolor{rowA}{RGB}{220,235,252}   % design closure + procurement
\definecolor{rowB}{RGB}{198,239,206}   % build
\definecolor{rowC}{RGB}{255,235,156}   % test + CAD/BOM
\definecolor{rowD}{RGB}{252,213,180}   % integrate + submit

% File-name macro: typewriter face with line breaks allowed after underscores,
% so long artifact names do not overflow the margin.
\newcommand{\f}[1]{{\ttfamily\def\_{\textunderscore\allowbreak}#1}}

\title{\textbf{Five-Day Project Wrap-Up Plan}\\[0.4em]
\large Autonomous Thermoelectric Atmospheric-Water PV-Cleaning Robot\\[0.2em]
\normalsize Design Closure, 4-Tile Bench Prototype, Test, and Thesis Submission\\
Thursday 11 June -- Monday 15 June 2026}
\author{Rahma El-Hamouly, Graduation Project}
\date{Issued: 10 June 2026}

\begin{document}
\maketitle

\begin{abstract}
\noindent
Five calendar days close the project: Day~1 resolves the eight open engineering
flags, freezes the design, and buys the parts; Days~2--3 build the 4-tile bench
demonstrator while thesis chapters that need no test data are assembled in
parallel; Day~4 runs the condensation test, the CAD block model, and the
full-scale BOM; Day~5 folds measured results into the thesis, compiles, runs
QA, and submits. The plan assumes one student working from Cairo with same-day
parts access (Ram Electronics, Future Electronics Egypt, local component
markets). The central technical risk is the test month: June is the worst
month for condensation in the annual simulation, so testing uses a humidified
enclosure indoors plus dawn-hour outdoor runs, and validates the condensation
model point-wise at measured air states rather than against the annual-mean
yield figure. Chapter mapping is in \f{THESIS\_OUTLINE.tex}.
\end{abstract}

\tableofcontents
\newpage

% =============================================================================
\section{Key Deadlines}
% =============================================================================
\begin{center}
\small
\renewcommand{\arraystretch}{1.35}
\begin{tabular}{@{}lll@{}}
\toprule
\textbf{Milestone} & \textbf{Date} & \textbf{Deliverable} \\
\midrule
\rowcolor{rowA}
Day 1: Design freeze + parts in hand & Thu 11 Jun 2026\,$^{*}$ & Freeze memo; all bench parts bought \\
\rowcolor{rowB}
Day 2: Bench mechanically complete   & Fri 12 Jun 2026 & TEC stack assembled; wiring done \\
\rowcolor{rowB}
Day 3: Firmware holds set-point      & Sat 13 Jun 2026 & Logged shakedown; run~1 started \\
\rowcolor{rowC}
Day 4: Measured yield; CAD + BOM     & Sun 14 Jun 2026 & Run-1 data reduced; CAD; 69-tile BOM \\
\rowcolor{rowD}
Day 5: \textbf{Submission}           & \textbf{Mon 15 Jun 2026} & Thesis PDF + packed folder \\
\bottomrule
\end{tabular}

{\footnotesize $^{*}$The freeze memo and the purchases must both land on the
11th. Friday shops open only after midday prayers, which is why all buying
happens Thursday.}
\end{center}

% =============================================================================
\section{Ground Rules}
% =============================================================================
\begin{enumerate}[leftmargin=*]
  \item \textbf{Build target is the 4-tile bench}, not the 69-tile system. The
        model is per-tile linear, so 4 tiles validate the physics and the
        full-scale numbers follow by multiplication
        (\f{prototype\_design\_report.tex}, prototype chapter).
  \item \textbf{The bench runs from a mains 12~V supply.} Battery, PV module,
        drive motors, and the cleaning head are not needed to validate
        condensation and control; they are bought only if budget and time
        allow (see the priority-cut list, \S\ref{sec:cuts}).
  \item \textbf{Design numbers change only on Day 1.} After the freeze memo,
        any change to TEC count, heat-sink resistance, battery size, or
        economics requires re-running \f{scripts/tec\_robot\_design.py} and
        noting the change in the memo. No silent edits.
  \item \textbf{Thesis assembly runs in the gaps.} Epoxy cure, paste settling,
        and overnight runs leave hours of dead time on Days 2--4; those hours
        go to the chapters that need no test data
        (\f{THESIS\_OUTLINE.tex} chs.~1--9, 11).
  \item \textbf{Every measurement is logged to CSV at the time it is taken},
        with a photo. Re-creating a lost data point costs a full night.
\end{enumerate}

% =============================================================================
\section{The June Testing Problem}
\label{sec:june}
% =============================================================================
The annual simulation puts June yield at essentially zero for Benban (night
dew points near 3--6~$^{\circ}$C at 17--20\% RH). Cairo in June is somewhat
more humid than Benban but outdoor midnight conditions are still marginal.
Two responses, both planned rather than improvised:

\begin{enumerate}[leftmargin=*]
  \item \textbf{Humidified-enclosure runs (primary).} The bench sits in a
        sealed clear storage box (60--80~L). An ultrasonic humidifier, or a
        tray of hot water refreshed from a kettle, raises the enclosure to
        roughly 28--32~$^{\circ}$C at 60--80\% RH, i.e.\ a dew point of
        20--26~$^{\circ}$C. The cold plate reaches 5--15~$^{\circ}$C under
        load, so the plate sits far below the dew point and condensation is
        guaranteed. The DHT22 measures the actual enclosure state, and the
        measured collection rate is compared with Eq.~(3.7) of
        \f{prototype\_design\_report.tex} evaluated at that same state. At
        30~$^{\circ}$C / 70\% RH the model predicts a collection rate of
        order tens of grams per hour for the 4-tile plate; even at one tenth
        of that, an overnight run gives a condensate mass a 0.1~g kitchen
        scale resolves to better than 5\%.
  \item \textbf{Dawn-hour outdoor run (secondary).} Cairo RH peaks between
        03:00 and 06:00. If a night with dew point above about
        12~$^{\circ}$C occurs during the window, one outdoor run at dawn
        provides a real-climate data point. This run is desirable, not
        required; the enclosure runs carry the validation.
\end{enumerate}

\textbf{What is being validated.} Not the 59~g/day figure directly. That
number is the 4-tile share of the \emph{Benban annual mean} (14.65~g/day per
tile) and June can never reproduce it outdoors. The thesis claim is: the
coupled heat/mass-transfer condensation model predicts the measured rate at
the measured air state within the uncertainty band. Once the model is
validated point-wise, the annual simulation built on it (and hence 59~g/day,
1.011~L/day, 69 tiles) stands. Chapter~10 of the thesis states this logic
explicitly so an examiner reads it as method, not excuse.

% =============================================================================
\section{Day 1 (Thursday 11 June): Design Closure and Procurement}
% =============================================================================

\subsection*{Morning (09:00--13:00): resolve the eight review flags}
Deliverable: a freeze memo, \f{docs/DESIGN\_FREEZE\_2026-06-11.md}, recording
one decision for each flag in \f{docs/REE598\_PROPOSAL\_REVIEW.md} \S3, plus
the regenerated numbers.

\begin{center}
\small
\renewcommand{\arraystretch}{1.25}
\begin{longtable}{@{}clp{9.2cm}@{}}
\toprule
\textbf{Flag} & \textbf{Topic} & \textbf{Decision to record} \\
\midrule
\endhead
3.1 & Tables 6.1/6.2 contradict & Regenerate Table 6.2 from the Table 6.1 soiling model. Scale B: 3{,}360 kWh/yr recovered, 67 USD/yr at 0.02 USD/kWh; payback recomputed at corrected unit cost. \\
3.2 & SEC optimism (0.927 kWh/L) & Supersede the proposal value with the coupled-solver SEC (4.16 kWh/L condensation-only, 5.27 with parasitics), inside the published 2.8--12 kWh/L band. \\
3.3 & Hot-side temperature & State that the coupled solver in \f{tec\_robot\_design.py} already iterates $T_H = T_{air} + Q_H R_{th}$; the proposal's decoupled COP is withdrawn. \\
3.4 & Heat-sink resistance & Re-run the simulation at $R_{th}$ = 0.15 / 0.5 / 1.0 K/W and keep the sensitivity table; spec a 100 mm+ fan-cooled extrusion (or CPU-cooler class sink) in the full-scale BOM. Bench measures the real value on Day 3. \\
3.5 & Scale-B battery undersized & BOM battery becomes a 2.5 kWh, 24 V LiFePO$_4$ pack with BMS. \\
3.6 & Battery price unrealistic & Pack priced at 450 USD; Scale-B unit total moves to about 900 USD; economics in flag 3.1 use this figure. \\
3.7 & Scale-C gantry infeasible & Scale C becomes a fleet of four or more Scale-B units; the gantry and its 1.67 t reservoir are dropped. \\
3.8 & Condensate tank too small & Operating sequence: harvest nightly into the 5 L on-board tank, pump out to a 100 L array-side buffer at dock; clean from the buffer. Sequence written into thesis \S4.5. \\
\bottomrule
\end{longtable}
\end{center}

Run \f{python scripts/tec\_robot\_design.py} for the $R_{th}$ sensitivity
(flag 3.4) before lunch so any change in tile count or SEC lands in the freeze
memo. Dependency: none. Fallback: if the sensitivity run stalls, freeze at the
existing design values and carry the conservative $R_{th}$ case as a
sensitivity table only; the flags are still all resolved on paper.

\subsection*{Afternoon (14:00--18:00): buy parts}
Deliverable: every bench-critical item physically in hand, receipts logged.
Route: Ram Electronics / Future Electronics Egypt for semiconductors, sensors,
and modules; local market for heat sinks, aluminium plate, fasteners, and the
storage box. Buy list and budget (from \f{prototype\_design\_report.tex}
BOM, plus the two bench additions):

\begin{center}
\small
\renewcommand{\arraystretch}{1.2}
\begin{longtable}{@{}lccr@{}}
\toprule
\textbf{Item} & \textbf{Qty} & \textbf{Unit (EGP)} & \textbf{Subtotal (EGP)} \\
\midrule
\endhead
TEC1-12706 Peltier module        & 4 & 130 & 520 \\
Cold-side heat sink + 40 mm fan  & 4 & 90  & 360 \\
Hot-side heat sink + 80 mm fan   & 4 & 150 & 600 \\
Anodised Al cold plate, 2 mm     & 1 & 120 & 120 \\
ESP32 dev board                  & 1 & 350 & 350 \\
DHT22 T/RH sensor (buy a spare)  & 2 & 180 & 360 \\
DS18B20 waterproof probe         & 2 & 60  & 120 \\
MOSFET PWM driver board          & 2 & 120 & 240 \\
Buck converter 12 V $\to$ var.   & 2 & 90  & 180 \\
INA219 current/voltage sensor    & 1 & 150 & 150 \\
12 V / 10 A mains PSU            & 1 & 400 & 400 \\
Thermal paste, clamps, frame, tubing, box, misc. & 1 & 700 & 700 \\
\midrule
\textbf{Bench-critical subtotal} & & & \textbf{4{,}100} \\
\midrule
\multicolumn{4}{@{}l}{\emph{Optional (only if time/budget allow; see cut list):}} \\
Micro diaphragm pump + tubing    & 1 & 250 & 250 \\
Brush + spray head               & 1 & 200 & 200 \\
DC gear motors + wheels          & 2 & 220 & 440 \\
12 V LiFePO$_4$ battery 7--12 Ah & 1 & 900 & 900 \\
50 Wp PV module                  & 1 & 700 & 700 \\
\midrule
\textbf{Full-kit total}          & & & \textbf{$\approx$6{,}590} \\
\bottomrule
\end{longtable}
\end{center}

The bench-critical subtotal ($\approx$4{,}100~EGP $\approx$ 85~USD at
48~EGP/USD) is what Days 2--4 actually need. Dependency: shops close around
21:00 but stock checks by phone happen in the morning, during flag work.
Fallback: any item missing at Ram/Future is ordered on Amazon.eg with next-day
delivery and the build order on Day~2 is rearranged so the missing item is
needed last; TEC1-12706, DHT22, ESP32, and 12706-class heat-sink/fan kits are
stocked by multiple Cairo vendors, so a full block is unlikely.

\subsection*{Evening (19:00--22:00): bench drawings and thesis skeleton}
Deliverables: (i) one-page bench layout sketch and wiring diagram
(ESP32 pins, MOSFET gate per TEC pair, sensor bus, INA219 shunt placement),
(ii) \f{REE599\_thesis.tex} skeleton created with all chapter files stubbed
per \f{THESIS\_OUTLINE.tex}, compiling to a PDF with front matter and empty
chapters. Dependency: freeze memo (the skeleton imports the corrected
numbers). Fallback: if the evening is lost to shopping overruns, the skeleton
moves to Day~2 gaps; the wiring sketch cannot slip because Day~2 morning
follows it.

% =============================================================================
\section{Day 2 (Friday 12 June): Mechanical Build + Electronics}
% =============================================================================
Friday is a home-workshop day by design; no shops needed.

\subsection*{Morning (09:00--13:00): TEC stack assembly}
Deliverable: four TEC sandwiches built: hot sink / paste / TEC / paste / cold
plate, clamped for full-face contact; cold plate tilted about 15$^{\circ}$
over a gutter and collection bottle; shared cold-side fan ducted across the
plate. Dependency: all Day-1 parts. Fallback: if clamping hardware
disappoints, zip-tie plus silicone interim assembly is acceptable for the
bench; note it in the build log.

\subsection*{Afternoon (14:00--18:00): wiring and power}
Deliverable: ESP32 wired to DHT22 (enclosure air), two DS18B20 (cold plate,
hot sink), INA219 in the TEC supply line, MOSFET PWM stage driving the TEC
pairs from the 12~V PSU; smoke test passed at low duty. Dependency: morning
assembly (sensor placement). Fallback: if the INA219 misbehaves, log voltage
and current manually with a multimeter at fixed intervals during runs; COP
accuracy drops but yield validation is untouched.

\subsection*{Evening (19:00--22:00): logging firmware + thesis chapters 1--3}
Deliverables: (i) firmware milestone~1: read all sensors at 5~s period, stream
CSV over serial, save to flash or laptop logger; (ii) thesis chapters 1--3
assembled from the sources mapped in \f{THESIS\_OUTLINE.tex} (introduction,
literature, Benban evidence), since none of them touch test data. Dependency:
afternoon wiring. Fallback: chapters slip to Day-3 gaps before firmware does;
the build has priority all evening if wiring debugging spills over.

% =============================================================================
\section{Day 3 (Saturday 13 June): Firmware, Shakedown, Run 1}
% =============================================================================

\subsection*{Morning (09:00--13:00): set-point hold and calibration}
Deliverables: (i) firmware milestone~2: PID (or fallback hysteresis) loop
holding a cold-plate set-point via PWM duty, anti-windup at the current limit;
(ii) calibration record: DS18B20 ice-point check (0.0~$^{\circ}$C), DHT22
cross-checked against a second sensor or local METAR. Dependency: Day-2
firmware milestone~1. Fallback: if PID tuning resists, ship bang-bang control
with a 1~K hysteresis band; it holds the plate below dew point just as well
for yield purposes and the thesis reports it honestly as the implemented
controller.

\subsection*{Afternoon (14:00--18:00): dry shakedown + enclosure build}
Deliverables: (i) one-hour powered run, no humidification: log plate
temperature vs.\ duty, hot-sink temperature, current; extract the measured
hot-side thermal resistance (closes review flag 3.4 with data); (ii) the
humidified enclosure (\S\ref{sec:june}) built and leak-checked, humidifier or
kettle-tray method tested to hold 60--80\% RH. Dependency: morning firmware.
Fallback: if the plate will not pull low enough, drop to two TECs per channel
in series-parallel or raise airflow; record whatever operating point is
reachable, since the model is validated at the measured state anyway.

\subsection*{Evening (19:00--23:00): start overnight run 1 + thesis chapters 4--7}
Deliverables: (i) run~1 started by 22:00 inside the enclosure: set-point
5~K below enclosure dew point, logging every 5~s, collection bottle weighed
dry beforehand; (ii) thesis chapters 4--7 assembled in the gaps (architecture,
TEC design, mechanics, power; all source-mapped, none test-dependent).
Dependency: afternoon enclosure. Fallback: if the enclosure cannot hold RH,
run anyway in the bathroom after a hot shower (crude but produces a 70\%+ RH
start) and fix the enclosure for run~2.

% =============================================================================
\section{Day 4 (Sunday 14 June): Test Data, CAD, Full-Scale BOM}
% =============================================================================

\subsection*{Morning (08:00--13:00): reduce run 1, set up run 2}
Deliverables: (i) condensate weighed (0.1~g scale), logs reduced to a
predicted-vs-measured table: measured $\dot m$ vs.\ the condensation equation
at the logged enclosure state, measured electrical power and estimated COP
from the INA219 trace, uncertainty per the propagation formula in the test
plan; (ii) run-2 decision: dawn outdoor run if tonight's forecast dew point
exceeds about 12~$^{\circ}$C, otherwise a second enclosure run at a different
set-point (model check at a second operating point). Dependency: run 1
completed. Fallback: if run 1 produced no measurable condensate (leak, plate
too warm), debug before noon and convert Day-4 night into the primary run;
the schedule holds because thesis integration is Day 5.

\subsection*{Afternoon (14:00--19:00): CAD block model + 69-tile BOM}
Deliverables: (i) Fusion~360 block-level assembly of the Scale-B robot
(chassis, TEC stack envelope, tanks, cleaning head, docking interface), one
exploded view and one dimensioned overall drawing, exported to
\f{figures/} for thesis chapter 6 and Appendix D; (ii) full-scale 69-tile BOM
with Egyptian vendor pricing (extend the bench BOM line items $\times$
quantities, add chassis, 2.5~kWh battery at the corrected price, 1{,}150~Wp
PV, MPPT), saved as Appendix B.2. Dependency: freeze memo numbers. Fallback:
if CAD time evaporates, ship the dimensioned block diagram from
\f{figures/system\_architecture.svg} annotated with the frozen dimensions, and
note the solid model as future work; the BOM cannot be cut.

\subsection*{Evening (19:00--23:00): chapter 10 draft + run 2}
Deliverables: (i) thesis chapter 10 drafted around the run-1 numbers
(build, calibration, protocol, June-method section, results table, uncertainty);
(ii) run~2 started (dawn outdoor or second enclosure point); (iii) chapters
8, 9, and 11 assembled (control design text exists; simulation chapter reads
off the CSV; economics rebuilt from the freeze memo). Dependency: morning data
reduction. Fallback: chapter 11 economics is the piece to thin out if hours
run short, down to the corrected Scale-B table plus one paragraph per scale.

% =============================================================================
\section{Day 5 (Monday 15 June): Integrate, Compile, QA, Submit}
% =============================================================================

\subsection*{Morning (08:00--13:00): close the data and the text}
Deliverables: run-2 condensate weighed and reduced; chapter 10 finalised with
both runs; chapter 8 \S\S8.5--8.6 written (firmware as built, measured
set-point hold); chapter 12 conclusions and the abstract written with final
numbers; firmware source and raw CSVs dropped into Appendices C and E.
Dependency: run 2. Fallback: one validated run is sufficient for the thesis
claim; if run 2 failed, chapter 10 reports run 1 plus the failure analysis,
which examiners read as competent experimental practice.

\subsection*{Afternoon (14:00--18:00): full compile and QA}
Deliverables: full build chain
(\f{pdflatex} $\to$ \f{bibtex} $\to$ \f{pdflatex} $\times$2) clean; QA pass
adapted from \f{submission/SUBMISSION\_CHECKLIST.md}:
\begin{itemize}[nosep,leftmargin=*]
  \item no \texttt{??} cross-references, no missing citations, all figures
        resolve;
  \item abstract numbers match \f{output/tec\_robot\_results.json} and the
        freeze memo;
  \item guideline format check: cover page, acknowledgment, abstract page,
        TOC/LOF/LOT, nomenclature, chapter-based numbering of
        sections/figures/tables/equations, 12~pt, 1.5 spacing;
  \item every review flag 3.1--3.8 traceable to a thesis section;
  \item spell-check pass.
\end{itemize}
Dependency: morning text closure. Fallback: QA items that find deep problems
(a wrong number propagated) take priority over polish items; fix numbers
first, typography last.

\subsection*{Evening (18:00--21:00): pack and submit}
Deliverable: submission folder packed and sent per department instructions:
\f{REE599\_thesis.pdf}, source \f{.tex} files, \f{references.bib},
\f{scripts/} and \f{output/} (simulation), firmware source, CAD exports, raw
test data CSVs, photos of the rig, the freeze memo, and the two planning
documents. Dependency: afternoon QA. Fallback: none needed; this block has
two hours of slack by design.

% =============================================================================
\section{Risk Table}
% =============================================================================
\begin{center}
\small
\renewcommand{\arraystretch}{1.25}
\begin{longtable}{@{}p{3.6cm}ccp{7.0cm}@{}}
\toprule
\textbf{Risk} & \textbf{Likelihood} & \textbf{Impact} & \textbf{Mitigation / fallback} \\
\midrule
\endhead
A part is out of stock Day 1 & Low & High & Phone stock check in the morning; two vendors per item; Amazon.eg next-day as backstop; build order rearranged so the late part is needed last. \\
No condensation in June Cairo air & High outdoors & High & Planned for, not residual: humidified enclosure is the \emph{primary} method (\S\ref{sec:june}); dew point inside is 20+~$^{\circ}$C, far above the plate. Dawn (03:00--06:00) outdoor run only as a bonus data point. \\
Cold plate will not get below dew point & Medium & High & Measured hot-side resistance from the Day-3 dry run; raise airflow, re-seat paste, or run fewer tiles harder; validation works at any reachable plate temperature. \\
DHT22 dead or drifting & Medium & Medium & Spare bought on Day 1; cross-check both sensors against each other and the ice-point-calibrated DS18B20 in a sealed bag. \\
INA219 unavailable or noisy & Medium & Low & Manual multimeter readings at 30-min marks; COP uncertainty grows, yield validation unaffected. \\
PID will not stabilise & Medium & Low & Hysteresis (bang-bang) controller with 1 K band; honest write-up. \\
ESP32 toolchain friction & Low & Medium & Fall back to Arduino Uno + SD shield (locally stocked); firmware is plain Arduino C either way. \\
Scale resolution too coarse & Low & Medium & Overnight integration grows the mass; if needed, a 0.01 g jewellery scale costs about 250 EGP at local markets. \\
Run 1 lost (leak, power cut) & Medium & Medium & Day-4 night is the built-in repeat; one good run suffices for the thesis claim. \\
LaTeX build breakage & Low & Medium & Thesis reuses the proven \f{prototype\_design\_report.tex} preamble and shim; skeleton compiles from Day 1 evening so breakage surfaces early. \\
Time overrun anywhere & Medium & High & Priority-cut list below; cuts are pre-decided so no time is lost deciding. \\
\bottomrule
\end{longtable}
\end{center}

% =============================================================================
\section{Priority-Cut List (If Behind Schedule)}
\label{sec:cuts}
% =============================================================================
Cut from the top; the bottom is never cut.
\begin{enumerate}[leftmargin=*]
  \item Mobility hardware (motors, wheels) and the cleaning-head demo: drop
        first; they are already covered by design calculations and the BOM.
  \item PV + battery charging of the bench: bench stays on the mains PSU; the
        power chapter is analytical anyway.
  \item Run 2: one validated run carries the claim.
  \item Fusion~360 solid model: replace with the annotated dimensioned block
        diagram; keep the drawing in Appendix D.
  \item INA219 instrumentation: manual meter logging.
  \item Chapter 11 depth: corrected Scale-B economics table plus one paragraph
        per scale.
  \item Chapter 3 length: compress further; the standalone Benban PDFs remain
        in the archive.
  \item \textbf{Never cut:} the eight flag resolutions and freeze memo; the
        overnight condensation run with a weighed yield; the 69-tile BOM;
        thesis chapters 1, 5, 9, 10, 12; the clean compile and QA pass.
\end{enumerate}

% =============================================================================
\section{Calendar Summary}
% =============================================================================
\begin{center}
\small
\renewcommand{\arraystretch}{1.3}
\begin{longtable}{@{}clp{8.6cm}@{}}
\toprule
\textbf{Day} & \textbf{Date} & \textbf{Blocks (morning / afternoon / evening)} \\
\midrule
\endhead
\rowcolor{rowA} 1 & Thu 11 Jun & Resolve 8 flags + freeze memo / buy all parts ($\approx$4{,}100 EGP bench-critical) / wiring sketch + thesis skeleton \\
\rowcolor{rowB} 2 & Fri 12 Jun & TEC stack assembly / ESP32 wiring + smoke test / logging firmware + thesis chs. 1--3 \\
\rowcolor{rowB} 3 & Sat 13 Jun & Set-point firmware + calibration / dry shakedown + enclosure / \textbf{overnight run 1} + thesis chs. 4--7 \\
\rowcolor{rowC} 4 & Sun 14 Jun & Reduce run 1 / CAD block model + 69-tile BOM / ch. 10 draft + \textbf{run 2} + chs. 8, 9, 11 \\
\rowcolor{rowD} 5 & Mon 15 Jun & Reduce run 2 + final chapters / full compile + QA / \textbf{pack and submit} \\
\bottomrule
\end{longtable}
\end{center}

% =============================================================================
\section{Build Commands (Project Root)}
% =============================================================================
\begin{verbatim}
# Regenerate simulation outputs after any frozen-parameter change (Day 1)
python scripts/tec_robot_design.py

# Compile the thesis (Days 1-5)
pdflatex REE599_thesis.tex
bibtex   REE599_thesis
pdflatex REE599_thesis.tex
pdflatex REE599_thesis.tex

# Compile the two planning documents (no bibliography)
pdflatex THESIS_OUTLINE.tex
pdflatex THESIS_OUTLINE.tex
pdflatex FIVE_DAY_WRAPUP_PLAN.tex
pdflatex FIVE_DAY_WRAPUP_PLAN.tex
\end{verbatim}

\end{document}
